\documentclass[10pt,a4paper]{amsart}
\usepackage{amsmath,amssymb,amsthm,graphicx}
% For commutative diagrams we use Xy-pic
\usepackage[all]{xy}

% We use 2cell for 2-commutative diagrams.
\xyoption{2cell}
% Theorem environments.
%
\theoremstyle{plain}
\newtheorem{theorem}[subsection]{Theorem}
\newtheorem{proposition}[subsection]{Proposition}
\newtheorem{lemma}[subsection]{Lemma}

\theoremstyle{definition}
\newtheorem{definition}[subsection]{Definition}
\newtheorem{example}[subsection]{Example}
\newtheorem{exercise}[subsection]{Exercise}
\newtheorem{situation}[subsection]{Situation}

\theoremstyle{remark}
\newtheorem{remark}[subsection]{Remark}
\newtheorem{remarks}[subsection]{Remarks}

\numberwithin{equation}{subsection}

% Macros
%
\def\lim{\mathop{\mathrm{lim}}\nolimits}
\def\colim{\mathop{\mathrm{colim}}\nolimits}
\def\Spec{\mathop{\mathrm{Spec}}}
\def\Hom{\mathop{\mathrm{Hom}}\nolimits}
\def\Ext{\mathop{\mathrm{Ext}}\nolimits}
\def\SheafHom{\mathop{\mathcal{H}\!\mathit{om}}\nolimits}
\def\SheafExt{\mathop{\mathcal{E}\!\mathit{xt}}\nolimits}
\def\Sch{\mathit{Sch}}
\def\Mor{\mathop{\mathrm{Mor}}\nolimits}
\def\Ob{\mathop{\mathrm{Ob}}\nolimits}
\def\Sh{\mathop{\mathit{Sh}}\nolimits}
\def\NL{\mathop{N\!L}\nolimits}
\def\CH{\mathop{\mathrm{CH}}\nolimits}
\def\proetale{{pro\text{-}\acute{e}tale}}
\def\etale{{\acute{e}tale}}
\def\QCoh{\mathit{QCoh}}
\def\Ker{\mathop{\mathrm{Ker}}}
\def\Im{\mathop{\mathrm{Im}}}
\def\Coker{\mathop{\mathrm{Coker}}}
\def\Coim{\mathop{\mathrm{Coim}}}

% Boxtimes
%
\DeclareMathSymbol{\boxtimes}{\mathbin}{AMSa}{"02}

%
% Macros for moduli stacks/spaces
%
\def\QCohstack{\mathcal{QC}\!\mathit{oh}}
\def\Cohstack{\mathcal{C}\!\mathit{oh}}
\def\Spacesstack{\mathcal{S}\!\mathit{paces}}
\def\Quotfunctor{\mathrm{Quot}}
\def\Hilbfunctor{\mathrm{Hilb}}
\def\Curvesstack{\mathcal{C}\!\mathit{urves}}
\def\Polarizedstack{\mathcal{P}\!\mathit{olarized}}
\def\Complexesstack{\mathcal{C}\!\mathit{omplexes}}
% \Pic is the operator that assigns to X its picard group, usage \Pic(X)
% \Picardstack_{X/B} denotes the Picard stack of X over B
% \Picardfunctor_{X/B} denotes the Picard functor of X over B
\def\Pic{\mathop{\mathrm{Pic}}\nolimits}
\def\Picardstack{\mathcal{P}\!\mathit{ic}}
\def\Picardfunctor{\mathrm{Pic}}
\def\Deformationcategory{\mathcal{D}\!\mathit{ef}}
\usepackage{cleveref}
\crefname{originallemma}{original lemma}{original lemmas}
\Crefname{originallemma}{Original lemma}{Original lemmas}
\makeatletter
\newlabel{descent-lemma-sheaf-condition-holds}{{03DT}{source}{}{sourcepointer-1}{}}
\newlabel{descent-lemma-sheaf-condition-holds@cref}{{[originallemma][0][]03DT}{[1][0][]source}{}{}{}}
\makeatother
\crefname{originallemma}{original lemma}{original lemmas}
\Crefname{originallemma}{Original lemma}{Original lemmas}
\makeatletter
\newlabel{lemma-fpqc-sheaves}{{03O1}{source}{}{sourcepointer-2}{}}
\newlabel{lemma-fpqc-sheaves@cref}{{[originallemma][0][]03O1}{[1][0][]source}{}{}{}}
\makeatother
\crefname{originallemma}{original lemma}{original lemmas}
\Crefname{originallemma}{Original lemma}{Original lemmas}
\makeatletter
\newlabel{lemma-descent-modules}{{03OA}{source}{}{sourcepointer-3}{}}
\newlabel{lemma-descent-modules@cref}{{[originallemma][0][]03OA}{[1][0][]source}{}{}{}}
\makeatother
\title{The presheaf of sections of a pullback of a quasi-coherent sheaf satisfies the fpqc sheaf condition}
\author{The Stacks Project Authors}
\date{stacks-project:etale-cohomology.tex:local-snapshot-20261011; GNU FDL 1.2 or later}
\begin{document}
\maketitle
\noindent\textit{Editorial note.} All delivered mathematical text is the author's own, copied from the pinned Stacks Project chapter file etale-cohomology.tex; no statement or proof step was rewritten, reordered or removed.\par\medskip
\noindent\textit{Editorial note.} The preamble is the chapter's own native preamble excerpt used for this source: the Xy-pic setup at lines 14-18 and the theorem declarations with the mathematical macros at lines 154-217 of sources/stacks-etale-cohomology/preamble.tex. The quoted class line and the xr-hyper externaldocument block are not included because the assembler supplies a plain amsart wrapper; no macro was invented or redefined.\par\medskip
\noindent\textit{Editorial note.} The two same-chapter pointers are located at their \textbackslash{}label line inside the pinned chapter file; the cross-chapter pointer is located in the project's own tag list staged in the snapshot (sources/stacks-chow/tags.txt), at the verbatim line carrying its canonical tag.\par\medskip
\noindent\textit{Allowed standard prerequisites (editorial).}\begin{itemize}
\item The fpqc topology and the sheaf condition for it. Standard in this chapter; the statement refers to the presheaf F\textasciicircum{}a on Sch/S and to the sheaf condition for the fpqc topology. Supplies the reduction of the sheaf condition to Zariski coverings and faithfully flat morphisms of affine schemes; the verification in both cases is the counted content.
\item Descent of modules along a faithfully flat ring map. Same chapter, Lemma labelled lemma-descent-modules (source line 1237, Stacks tag 03OA); invoked at proof lines 1364-1365. Supplies exactness of 0 -> M -> B tensor\_A M -> B tensor\_A B tensor\_A M for A -> B faithfully flat; the translation of that exactness into the sheaf condition is retained in the counted span.
\end{itemize}\medskip
\ifdefined\hypertarget\hypertarget{sourcepointer-1}{}\fi
\noindent\textit{Source reference.} Original lemma 03DT: sources/stacks-chow/tags.txt, source lines 4432--4432. Cross-chapter pointer named at proof line 1346 as where the proposition is proved; the author immediately announces 'We indicate the proof here' and the indicated argument in this chapter is retained verbatim in the counted span.\par\medskip
\ifdefined\hypertarget\hypertarget{sourcepointer-2}{}\fi
\noindent\textit{Source reference.} Original lemma 03O1: sources/stacks-etale-cohomology/pinned-original-source.tex, source lines 974--974. Same-chapter pointer invoked at proof lines 1347-1349 for the reduction of the fpqc sheaf condition to Zariski coverings and faithfully flat morphisms of affine schemes; both cases of the verification are retained in the counted span.\par\medskip
\ifdefined\hypertarget\hypertarget{sourcepointer-3}{}\fi
\noindent\textit{Source reference.} Original lemma 03OA: sources/stacks-etale-cohomology/pinned-original-source.tex, source lines 1237--1237. Same-chapter pointer invoked at proof lines 1364-1365 for exactness of the Amitsur complex of a faithfully flat ring map; the identification of that exactness with the sheaf condition for F\textasciicircum{}a is retained in the counted span.\par\medskip
\begin{proposition}
\label{proposition-quasi-coherent-sheaf-fpqc}
For any quasi-coherent sheaf $\mathcal{F}$ on $S$ the presheaf
$$
\begin{matrix}
\mathcal{F}^a : & \Sch/S & \to & \textit{Ab}\\
& (f: T \to S) & \mapsto & \Gamma(T, f^*\mathcal{F})
\end{matrix}
$$
is an $\mathcal{O}$-module which satisfies the sheaf condition for the
fpqc topology.
\end{proposition}

\begin{proof}
This is proved in
Descent, Lemma \ref{descent-lemma-sheaf-condition-holds}.
We indicate the proof here. As established in
Lemma \ref{lemma-fpqc-sheaves},
it is enough to check the sheaf property
on Zariski coverings and faithfully flat morphisms of affine schemes. The
sheaf property for Zariski coverings is standard scheme theory, since
$\Gamma(U, i^\ast \mathcal{F}) = \mathcal{F}(U)$ when
$i : U \hookrightarrow S$ is an open immersion.

\medskip\noindent
For $\left\{\Spec(B)\to \Spec(A)\right\}$ with $A\to B$ faithfully
flat and
$\mathcal{F}|_{\Spec(A)} = \widetilde{M}$
this corresponds to the fact that
$M = H^0\left((B/A)_\bullet \otimes_A M \right)$, i.e., that
\begin{align*}
0 \to M \to B \otimes_A M \to B \otimes_A B \otimes_A M
\end{align*}
is exact by
Lemma \ref{lemma-descent-modules}.
\end{proof}

\end{document}
